\documentclass[10pt,a4paper]{amsart}
\usepackage[margin=19mm]{geometry}
\usepackage{amsmath,amssymb,mathtools}
\usepackage[scr=rsfs,cal=euler]{mathalfa}
\usepackage{xfrac}
\usepackage[unicode,hidelinks,bookmarks=false]{hyperref}
\newcommand{\ie}{i.e.\ }
\newcommand{\cf}{cf.\ }
\newcommand{\resp}{respectively\ }
\newcommand{\Subsection}{\subsection}
\providecommand{\nobreakdash}{\nobreak\hbox{-}}
\setlength{\emergencystretch}{2em}
\multlinegap0pt
\usepackage{amssymb}
\usepackage[normalem]{ulem}
\usepackage{bm}
\usepackage{mathtools}
\newtheorem{lem}{Lemma}
\title{Equilibrium measures for higher dimensional \\ rotationally symmetric Riesz gases}
\author{Sung-Soo Byun, Peter J. Forrester, Satya N. Majumdar, Gregory Schehr}
\date{Source: arXiv:2602.03047v1 (CC BY 4.0)}
\begin{document}
\maketitle

\noindent\emph{Source and licence.} Equilibrium measures for higher dimensional \\ rotationally symmetric Riesz gases. Version arXiv:2602.03047v1 (CC BY 4.0), distributed by arXiv under the Creative Commons Attribution 4.0 International licence, http://creativecommons.org/licenses/by/4.0/, as recorded in the arXiv licence metadata for this exact version and shown on the abstract page https://arxiv.org/abs/2602.03047v1. Full licence notice: \texttt{licenses/arxiv-2602.03047-CC-BY-4.0.txt}.\par\smallskip

\noindent\textit{Editorial scope.} Counted unit: the article's lem potential (source0.tex lines 1689--1694). The article's complete original proof of it is delivered with it (source0.tex lines 1695--1750, one \texttt{proof} environment of 56 lines). No mathematics is written, repaired or re-derived: the statement and the proof are the article's own lines, in the article's own notation, and no retained line is substituted or re-flowed.  Retained uncounted as the source-local context the proof needs: lines 570--572.  Nothing was omitted from any retained span, so the package carries no omission record.  No retained line contains a citation, so the wrapper carries no bibliography and no bibliography entry.  No retained line contains a figure environment, an image-inclusion command or drawing code, so no image is delivered and the package declares no asset dependency.  Editorial formatting only, in addition: an AMS \texttt{amsart} preamble that restates the article's own declarations and macros used by the retained lines, namely the environments \texttt{lem} on separate counters, together with the standard packages those lines need.  The retained environments print in this document as Lemma 1 in order of appearance, whereas the article prints the same items with the numbering of its own full document.  The complete native source of the article as distributed in the arXiv e-print for this exact version is bundled unchanged as \texttt{sources/193880/source0.tex}.
\begin{equation} \label{eq for binomial expansion}
(1-z)^{\alpha}= \sum_{k=0}^\infty \frac{ \Gamma(-\alpha+k) }{ \Gamma(-\alpha) k! } z^k, \qquad (|z|<1). 
\end{equation}

\begin{lem}\label{Lem_hypergeometric inequal for power type potential}
Let $d \ge 1$ and $s \in [d-2,d)$. Then for any $p\in \mathbb{Z}_{ \ge 0 }$ and $x \ge 1$, we have 
\begin{align} \label{2.34}
 \frac{\sin(\pi\tfrac{d-s}{2})}{\pi} \sum_{n=0}^\infty \frac{  \Gamma( \frac{s-d}{2}+n+1 )  }{  (n+\frac{s}{2}) n! }   \frac{2p}{2n+2p+s}    \, \frac{1}{x^{2n+s}}  \ge     \frac{ \Gamma(\frac{s}{2})   }{ \Gamma(\frac{d}{2})  }  \,  -    \frac{ \Gamma(\frac{s}{2}+p)   }{ \Gamma(\frac{d}{2}+p)  } \,      x^{2p}.
\end{align}
\end{lem}

\begin{proof} 

We first consider the case $s >0$. Then by \eqref{eq for binomial expansion}, we have 
\begin{align*}
 \sum_{n=0}^\infty \frac{  \Gamma( \frac{s-d}{2}+n+1 )  }{  (n+\frac{s}{2}) n! }   \frac{2p}{2n+2p+s} y^{n} &= 2\sum_{n=0}^\infty \frac{  \Gamma( \frac{s-d}{2}+n+1 )  }{    n! }  \bigg(\int_0^1 u^{2n+s-1}(1-u^{2p})\,du \bigg)  y^n
\\
&= 2 \Gamma( 1+\tfrac{s-d}{2} ) \int_0^1 u^{s-1} (1-u^2y)^{ \frac{d-s}{2}-1 } (1-u^{2p})\,du .
\end{align*}
Here, the condition $s>0$ is required for the integrability when $n=0$.
Then, the LHS of \eqref{2.34} is written as 
\begin{align}
\begin{split}
 \frac{2}{ \Gamma( \frac{d-s}{2} ) } \frac{1}{x^s} \int_0^1 u^{s-1} (1-u^2/x^2)^{ \frac{d-s}{2}-1 } (1-u^{2p})\,du
= \frac{2}{ \Gamma( \frac{d-s}{2} ) }   \int_0^{1/x} t^{s-1} (1-t^2)^{ \frac{d-s}{2}-1 } (1-t^{2p}x^{2p})\,dt. 
\end{split}  
\end{align}
Note that for $t>1/x$, we have $1-t^{2p}x^{2p}<0$. Therefore 
\begin{align*}
&\quad \frac{2}{ \Gamma( \frac{d-s}{2} ) }   \int_0^{1/x} t^{s-1} (1-t^2)^{ \frac{d-s}{2}-1 } (1-t^{2p}x^{2p})\,dt 
\\
& \ge \frac{2}{ \Gamma( \frac{d-s}{2} ) }   \int_0^{1} t^{s-1} (1-t^2)^{ \frac{d-s}{2}-1 } (1-t^{2p}x^{2p})\,dt =   \frac{ \Gamma(\frac{s}{2})   }{ \Gamma(\frac{d}{2})  }  \,  -    \frac{ \Gamma(\frac{s}{2}+p)   }{ \Gamma(\frac{d}{2}+p)  } \,      x^{2p},
\end{align*}
where the equality follows from the evaluation of the Euler beta integral in terms of gamma functions.
This completes the proof for the case $s>0$. 

For the remaining case $s \in (-1,0]$ when $d=1$, a slight modification is required to ensure integrability. First, note that 
\begin{align*}
&\quad \sum_{n=0}^\infty \frac{  \Gamma( \frac{s-d}{2}+n+1 )  }{  (n+\frac{s}{2}) n! }   \frac{2p}{2n+2p+s} y^{n} = \frac{  \Gamma( \frac{s-d}{2}+1 )  }{  s }   \frac{4p}{2p+s}+  \sum_{n=1}^\infty \frac{  \Gamma( \frac{s-d}{2}+n+1 )  }{  (n+\frac{s}{2}) n! }   \frac{2p}{2n+2p+s} y^{n} 
\\
&= \frac{  \Gamma( \frac{s-d}{2}+1 )  }{  s }   \frac{4p}{2p+s}+  2\sum_{n=1}^\infty \frac{  \Gamma( \frac{s-d}{2}+n+1 )  }{    n! }  \bigg(\int_0^1 u^{2n+s-1}(1-u^{2p})\,du \bigg)  y^n
\\
&=\frac{  \Gamma( \frac{s-d}{2}+1 )  }{  s }   \frac{4p}{2p+s}+ 2 \Gamma( \tfrac{s-d}{2}+1 ) \int_0^1 u^{s-1} \Big( (1-u^2y)^{ \frac{d-s}{2}-1 }-1\Big) (1-u^{2p})\,du .
\end{align*} 
Then the LHS of \eqref{2.34} is written as 
\begin{align}
\begin{split}
& \quad  \frac{2}{ \Gamma( \frac{d-s}{2} ) }  \frac{2p}{s(2p+s)} \frac{1}{x^s} + \frac{2}{ \Gamma( \frac{d-s}{2} ) } \frac{1}{x^s} \int_0^1 u^{s-1} \Big( (1-u^2/x^2)^{ \frac{d-s}{2}-1 } -1 \Big) (1-u^{2p})\,du
\\
&=  \frac{2}{ \Gamma( \frac{d-s}{2} ) }  \frac{2p}{s(2p+s)} \frac{1}{x^s} +  \frac{2}{ \Gamma( \frac{d-s}{2} ) }   \int_0^{1/x} t^{s-1} \Big( (1-t^2)^{ \frac{d-s}{2}-1 }-1\Big) (1-t^{2p}x^{2p})\,dt. 
\end{split}  
\end{align}
Here, again by using the Euler beta integral, we obtain  
\begin{align*}
&\quad \frac{2}{ \Gamma( \frac{d-s}{2} ) }   \int_0^{1/x} t^{s-1} \Big( (1-t^2)^{ \frac{d-s}{2}-1 }-1\Big) (1-t^{2p}x^{2p})\,dt 
\\
&\ge \frac{2}{ \Gamma( \frac{d-s}{2} ) }   \int_0^{1} t^{s-1} \Big( (1-t^2)^{ \frac{d-s}{2}-1 }-1\Big) (1-t^{2p}x^{2p})\,dt 
= \frac{2}{ \Gamma( \frac{d-s}{2} ) }\Big( \frac{x^{2p}}{s+2p} -\frac{1}{s}\Big)   + \frac{ \Gamma(\frac{s}{2})   }{ \Gamma(\frac{d}{2})  }   -    \frac{ \Gamma(\frac{s}{2}+p)   }{ \Gamma(\frac{d}{2}+p)  } \,      x^{2p}.
\end{align*}
Therefore, we have shown that 
\begin{align*}
&\quad  \frac{\sin(\pi\tfrac{d-s}{2})}{\pi} \sum_{n=0}^\infty \frac{  \Gamma( \frac{s-d}{2}+n+1 )  }{  (n+\frac{s}{2}) n! }   \frac{2p}{2n+2p+s}    \, \frac{1}{x^{2n+s}}
\\
& \ge \frac{2}{ \Gamma( \frac{d-s}{2} ) } \frac{1}{s(2p+s)} \Big( 2p(x^{-s}-1)+s (x^{2p}-1) \Big)  + \frac{ \Gamma(\frac{s}{2})   }{ \Gamma(\frac{d}{2})  }   -    \frac{ \Gamma(\frac{s}{2}+p)   }{ \Gamma(\frac{d}{2}+p)  } \,      x^{2p}.
\end{align*}
Since $s<0$ and $p>0$, the first term in the RHS is nonnegative for $x \ge 1$. Therefore the proof is complete. 
\end{proof}

\end{document}
